\chapter{Discussion} \label{chapter:ch6_discussion}

The introduction asked whether frequently changing collaborative-robot tasks can be programmed more directly by specifying spatial intent in the real workspace instead of relying on hand-guiding or editing motion commands and tool actions in a teach pendant or offline programming workflow. Within the scope of this prototype, the result is positive for selected task-oriented workflows.

This chapter interprets that result rather than recapping the implementation. The thesis does not propose a general robot programming environment. It shows that a selected use case can combine workspace sensing, spatial input, task-specific interpretation, and robot execution into one authoring workflow, and it also makes clear where that result currently stops.

\section{Main Result and Thesis Closure}

The main result is an end-to-end prototype for in-workspace spatial authoring on a collaborative robot. The implemented system combines a tracked 6DoF input device, stereo-vision-based sensing, task-specific interpretation, and robot execution behind a shared runtime with explicit interfaces. In practical terms, the prototype covers the full path from setup and calibration through authoring, confirmation, visualization, execution, and persistence.

At prototype scope, this closes the three goals stated in the introduction. For \goalref{goal:1}, the thesis demonstrates that poses and trajectories can be authored directly in the workspace and turned into executable robot behavior through task-specific interpretation. For \goalref{goal:2}, it provides a modular system in which different use cases reuse the same sensing, orchestration, and execution infrastructure. For \goalref{goal:3}, it demonstrates that approach on a real robot with two representative task families, pick-and-place and seam welding, while keeping the evaluation qualitative.

What matters here is not the shared runtime by itself, but the reusable operator workflow it supports. Across both use cases, the operator can select or configure a use case, capture the required spatial and non-spatial inputs, confirm the interpreted result, execute it, and save or load the authored program.

That common path is why the thesis did not produce only two isolated demonstrations. It produced a prototype in which different task-oriented procedures can share the same acquisition, confirmation, execution, and persistence flow.

That result is still intentionally bounded. The system does not expose a full robot programming language, and it does not try to cover rich branching, loops, exception recovery, or complex sensor-driven logic. Its stronger result is narrower: some frequently changing tasks can be moved from repeated low-level program editing to re-authoring inside a use case that already defines most of the task structure.

\section{Task Fit and Use-Case Meaning}

The two demonstrated task families are not equivalent examples of the same success. They support different conclusions about where the approach is useful and what kind of task structure it fits best.

Pick-and-place is the clearest evidence that the prototype can combine authored intent with refreshed scene information at execution time. The authored command does not store only a fixed world-frame pick pose. It stores an object identity and a pick pose relative to that object, then reconstructs the world pick pose from a fresh scene estimate when the command is executed. This shows that the system is not limited to replaying fixed authored motions. At the same time, that result is bounded by the current scene model and by the narrow handling task that was demonstrated.

Seam welding is the clearest practical fit for the present prototype. The authored trajectory expresses the main task intent directly in the workspace, and the use case can infer the surrounding robot procedure from that trajectory together with the detected scene geometry. In the demonstrated workflow, the user scans the scene, draws the weld, confirms the interpreted result, and then executes an approach--weld--depart sequence. That is a good match for the thesis question because it removes much of the low-level motion specification while keeping the authored input concrete and easy to interpret.

The usefulness of that workflow depends strongly on task context. For stable repetitive production, existing programming workflows may already be sufficient once the task has been commissioned. The present approach becomes more interesting when the task family stays similar but the concrete welds change often enough that repeated re-teaching becomes a noticeable burden, for example in small-batch or frequently re-taught work. In that setting, the human provides task intent in context and the robot performs the repeated physical execution.

More broadly, the thesis supports selected use cases with explicit task definitions, not general robot programming. A narrow use case can expose only the spatial inputs that actually vary and keep authoring direct. A broader use case can cover more variation, but it must also expose more parameters and more internal choices. The thesis does not resolve the best balance between those two directions, but it does show that both can be built on the same runtime structure as long as the task remains inside a prepared use case rather than a full programming environment.

\section{Applicability Boundary and Limitations}

The first major boundary is scene understanding. The current system detects only predefined tagged objects. Each scene read is based on one stereo image pair and returns only a small scene estimate, mainly object identities and poses. That is sufficient for the demonstrated pick-and-place and weld interpretation logic, but it is also the main reason the result should not be generalized too far beyond those examples. Richer tasks would require broader scene semantics, stronger robustness, and better handling of uncertainty than the current prototype provides.

The next boundary is sensing and calibration robustness. Pen tracking depends on marker visibility, lighting conditions, motion blur, and corner detection quality, and the current pen implementation does not fuse the available IMU data. Camera calibration also assumes that the stereo rig remains rigid with respect to the robot flange. These are workable assumptions for the thesis prototype, but they limit how confidently the present sensing stack can be transferred to less controlled conditions.

Execution remains equally bounded. The thesis demonstrates a shared robot motion interface and shows that use cases can drive real robot execution through it, but generalized tool control is not yet standardized and the evaluated use cases represent gripper and welding actions with fixed delays. No collision checking or simulated execution is provided.

The current use cases also rely on strong task-specific assumptions. The weld command assumes that the workpiece configuration stays unchanged after confirmation. The pick-and-place command assumes that the scene can be refreshed at execution time and that the expected object can still be matched by its stored identity.

UI maturity and portability are secondary limits. The frontend already covers setup, authoring, execution, and persistence through one workflow, yet some behaviors still reflect prototype maturity. Visualization can require manual reload, saved setups may restore only partially if hardware modules cannot reactivate, camera acquisition is currently Windows-only, and the current BLE discovery path differs across Windows and Linux.

Taken together, these limits define the applicability boundary of the work. The thesis supports selected task-oriented workflows in which a use case captures most of the task structure and the operator redefines the concrete spatial intent in the real workspace. It does not establish a general authoring method for arbitrary robot tasks, and it does not yet provide the scene understanding, execution safety, or tooling maturity required for broader deployment claims.

The evaluation boundary is equally explicit. The thesis demonstrates qualitative feasibility with limited physical validation, and most development and testing took place in the REMII offline environment. It does not include a formal user study, a controlled comparison against teach pendant or offline programming workflows, or quantitative characterization of tracking accuracy, repeatability, or productivity. The result should therefore be read as a proof-of-concept for selected workflows, not as evidence of measured superiority or industrial readiness.
